% ============================================================================
%  Exercises --- Intro to Pentesting, Red Teaming, Kali & Metasploit
%  Compile with:  pdflatex exercises.tex
% ============================================================================
\documentclass[11pt,a4paper]{article}

\usepackage[utf8]{inputenc}
\usepackage[T1]{fontenc}
\usepackage[margin=2.4cm]{geometry}
\usepackage{lmodern}
\usepackage{enumitem}
\usepackage{listings}
\usepackage{xcolor}
\usepackage{titlesec}
\usepackage{fancyhdr}
\usepackage{tcolorbox}
\usepackage{hyperref}
\hypersetup{colorlinks=true, linkcolor=blue!50!black, urlcolor=blue!50!black}

\definecolor{codeBg}{HTML}{F4F6F8}
\definecolor{kaliBlue}{HTML}{367BF0}
\definecolor{warnRed}{HTML}{C0392B}

\lstset{
  basicstyle=\ttfamily\small,
  backgroundcolor=\color{codeBg},
  breaklines=true,
  frame=single,
  rulecolor=\color{gray!40},
  showstringspaces=false,
}

\newtcolorbox{lab}{colback=kaliBlue!6, colframe=kaliBlue, sharp corners,
  boxrule=0.6pt, left=6pt, right=6pt, top=4pt, bottom=4pt}
\newtcolorbox{warn}{colback=warnRed!6, colframe=warnRed, sharp corners,
  boxrule=0.6pt, left=6pt, right=6pt, top=4pt, bottom=4pt}

\titleformat{\section}{\large\bfseries\color{kaliBlue}}{\thesection.}{0.5em}{}

\pagestyle{fancy}\fancyhf{}
\lhead{Offensive Security --- Exercises}
\rhead{\thepage}
\setlength{\headheight}{14pt}

\newcounter{taskc}[section]
\newcommand{\task}[1]{\refstepcounter{taskc}\medskip\noindent%
  \textbf{Exercise \thesection.\thetaskc \quad}\textit{(#1)}\par\nobreak\smallskip}
\newcommand{\pts}[1]{\hfill{\small\color{gray}[#1 pts]}}

\begin{document}

\begin{center}
  {\LARGE\bfseries Offensive Security --- Exercises}\\[0.4em]
  {\large Penetration Testing, Red Teaming, Kali Linux \& Metasploit}\\[0.4em]
  Cyber Security Master Lecture \hfill \today
\end{center}
\vspace{0.5em}

\begin{warn}
\textbf{Authorization \& ethics.} Perform the hands-on tasks \emph{only} inside
the isolated lab described below, or on platforms that explicitly permit it
(TryHackMe, Hack The Box, your own VMs). Scanning, exploiting, or even
probing systems you do not own --- without written permission --- is a criminal
offence (e.g.\ \S202a--c StGB in Germany). By doing these exercises you confirm
you act only against your own lab.
\end{warn}

\bigskip
\section*{Lab setup (do this first)}
\begin{lab}
\textbf{Goal:} a fully isolated network with one attacker and one target.
\begin{enumerate}[leftmargin=1.4em, itemsep=2pt]
  \item Install VirtualBox or VMware (free).
  \item Import the \textbf{Kali Linux} VM image
        (\url{https://www.kali.org/get-kali/}).
  \item Import \textbf{Metasploitable 2}
        (\url{https://sourceforge.net/projects/metasploitable/}).
  \item Set \emph{both} VMs' network adapter to \textbf{Host-Only} (or an
        Internal Network). \textbf{No bridged/NAT adapter on Metasploitable.}
  \item Boot both, log in to Metasploitable (\texttt{msfadmin}/\texttt{msfadmin}),
        run \texttt{ip a} to note its IP. From Kali, \texttt{ping} it to confirm
        connectivity \emph{and} that it has no internet route.
\end{enumerate}
Throughout, \texttt{TARGET} = the Metasploitable IP, \texttt{LHOST} = your Kali IP.
\end{lab}

% ---------------------------------------------------------------------------
\section{Concepts \& Ethics (no computer needed)}

\task{terminology} Define and contrast \emph{vulnerability scan},
\emph{penetration test}, and \emph{red team engagement}. For each, state the
primary goal and a typical deliverable. \pts{6}

\task{boxes} A client gives your team a low-privilege VPN account and an
internal network diagram, but no source code. Which testing model is this
(black/grey/white box)? Name one advantage and one limitation versus a pure
black-box test. \pts{4}

\task{legal} You are asked to ``just have a quick look'' at a partner company's
public website for vulnerabilities --- no contract yet, ``it's fine, we know
them''. List three documents/conditions that must exist before you run a single
\texttt{nmap} command, and explain the legal risk of proceeding without them.
\pts{6}

\task{killchain} Map the following actions to the correct phase of the
attack lifecycle (Recon, Scanning/Enum, Exploitation, Post-Exploitation,
Lateral Movement, Reporting): \emph{(a)} dumping local password hashes,
\emph{(b)} querying certificate-transparency logs for subdomains,
\emph{(c)} using stolen hashes to authenticate to a second server,
\emph{(d)} writing CVSS scores for each finding. \pts{4}

% ---------------------------------------------------------------------------
\section{Reconnaissance \& Enumeration (lab)}

\task{nmap-basics} From Kali, run a service/version scan of all ports on the
target and save the output:
\begin{lstlisting}[language=bash]
nmap -sV -p- -oA recon_target <TARGET>
\end{lstlisting}
List five open services with their version strings. Which \emph{one} immediately
looks suspicious/backdoored, and why? \pts{6}

\task{nmap-scripts} Re-run \texttt{nmap} with the default script set
(\texttt{-sC}) against the same host. Pick two scripts that fired and explain in
one sentence each what information they revealed. \pts{4}

\task{passive} Without touching the target, explain how you would discover
an organisation's externally-facing assets using only \emph{passive} techniques.
Name three distinct sources (e.g.\ DNS, certificate transparency, Shodan) and
what each yields. \pts{6}

% ---------------------------------------------------------------------------
\section{Exploitation with Metasploit (lab)}

\task{msf-workflow} Start \texttt{msfconsole}. Using the vsftpd finding from
Exercise~2.1, search for and select the matching exploit module, set the
required option(s), and run it. Paste the commands you used and the output
proving you obtained a shell. What privilege level did you land with
(\texttt{whoami})? \pts{8}

\task{module-types} For the module you just used, identify whether it is an
\texttt{exploit}, \texttt{auxiliary}, or \texttt{post} module. Then list one
example of each of the other two module types available in Metasploit and what
they are used for. \pts{6}

\task{payload-choice} Explain the difference between a \emph{bind} and a
\emph{reverse} payload. In a typical engagement where the target sits behind a
firewall that blocks inbound connections but allows outbound, which one do you
choose and why? \pts{4}

\task{msfvenom} Using \texttt{msfvenom}, generate a standalone Linux reverse-shell
payload (you do \emph{not} need to deliver it) pointing back to your Kali host on
port 4444. Provide the full command. Then describe how you would set up the
matching \texttt{multi/handler} listener in \texttt{msfconsole} to catch the
connection. \pts{6}

% ---------------------------------------------------------------------------
\section{Post-Exploitation \& Red Team Thinking}

\task{postex} You have a shell on the target. List, in order, the first four
things you would attempt and the goal of each (e.g.\ situational awareness,
privilege escalation, credential harvesting, persistence). \pts{6}

\task{pivot} Briefly explain what \emph{pivoting} / lateral movement means and
why a single compromised host loops you back to the ``scanning \& enumeration''
phase. Give one Metasploit/Meterpreter feature that enables pivoting. \pts{4}

\task{redvsblue} A red team gains domain admin but is detected by the SOC after
two days. Was the engagement a success or a failure? Argue both sides, then
state what the \emph{primary} deliverable of a red team engagement is (versus a
pentest). \pts{6}

\task{defence} For \emph{three} techniques you used or read about today
(e.g.\ the vsftpd backdoor exploit, hash dumping, reverse shells), name a
concrete defensive control that would prevent or detect it. \pts{6}

% ---------------------------------------------------------------------------
\section*{Challenge (optional, bonus)}
\begin{lab}
\task{capstone} Treat the lab as a mini-engagement. Produce a one-page
``finding'' for the most critical vulnerability you exploited, in proper report
format: title, severity (assign a CVSS vector), affected asset, description,
reproduction steps, evidence, and remediation. \pts{10 bonus}
\end{lab}

\vfill
\begin{center}\small
Total: 100 points (+10 bonus). Solutions are provided in a separate file.
\end{center}

\end{document}
